\documentclass{ximera}


% MATH -----------------------------------------------------------
\newcommand{\nats}{\mathbb N}
\newcommand{\ints}{\mathbb Z}
\newcommand{\rats}{\mathbb Q}
\newcommand{\reals}{\mathbb R}
\newcommand{\complex}{\mathbb C}
\newcommand{\powerset}{\mathscr P}

%-------------------------------------------------------------

\pgfplotsset{compat=1.18}
\begin{document}
\begin{abstract}
 \end{abstract}

    \title{Class Discussion Activity 3}
    
    \maketitle

\section*{Exercises}

\begin{enumerate}[label=(\arabic*)]



\item  Suppose $y(t)$ satisfies the Bernoulli IVP $y\,'=\dfrac{y}{x}+\dfrac{1}{y}$, $y(1)=1$.  Find $y(2026)$.  Round to the nearest hundredth.

% y_1=x:  y_1'-x/y_1=1-x/x=0.

% y=ux.  y^{\,\prime}=u'x+u.   u'x=1/(ux) ---> uu'=1/x^2 ---> u=\pm \sqrt(C-2/x).  y= x \sqrt(3-2/x)

% SANITY CHECK: y^{\,\prime}= \sqrt(3-2/x)+(1/x) 1/sqrt(3-2/x) = y/x+1/y.  YUP!

% y(2026)\approx 2026*sqrt(3-2/2026)\approx 3508.56


\item Suppose $y(t)$ satisfies the Bernoulli IVP $y\,'=\dfrac{1}{20}y-\dfrac{1}{50}y^2$, $y(0)=0$.  Find $y(2026)$.  Round to the nearest tenth.

% 0


% ---------------------------------------------------
% OLD EXERCISE:

% \item Suppose $y(t)$ satisfies the Bernoulli IVP $y\,'=\dfrac{1}{20}y-\dfrac{1}{50}y^2$, $y(0)=2$.  Find $y(2024)$ to the nearest tenth.

% 2.5, duh!

% Hard way:  100y^{\,\prime}=5y-2y^2 ---> 100y^{\,\prime}-5y=-2y^2.  y_1'-y_1/20=0 ---> y_1=e^(t/20).

% y=ue^(t/20).  y^{\,\prime}=u'e^(t/20)+(1/20)ue^(t/20).

% 100(u'e^(t/20)+(1/20)ue^(t/20))-5ue^(t/20)=-2(ue^(t/20))^2

% 100u'e^(t/20)=-2u^2 e^(2t/20)

% 100/u^2 u'=-2e^(t/20)

% -100/u =C-40e^(t/20)

% u/100 =1/(C+40e^(t/20))

% y = 100e^(t/20)/(10+40e^(t/20))=10/(e^(-t/20)+4)\approx 2.5 (for any reasonably large number)

%-------------------------------------------------------------------------------------



 \item Which of the following equations are Bernoulli equations?   Select all that apply.

\begin{enumerate}[label=(\alph*)]
\item $y^{\,\prime}=xy$
\item $y^{\,\prime}=xy+y^2$ % ***
\item $y^{\,\prime}=x^2y+x$
\item $y^{\,\prime}=x^2y+y$
\item $y^{\,\prime}=y+x\sqrt{y}$  % ***
\item $y^{\,\prime}=x^2y+y^2+x$ 
\item $y^{\,\prime}=x^2y+y+\frac{1}{y}$   % ***
\item $y^{\,\prime}=(y^2+1)e^x$  
\item $y^{\,\prime}=e^x y +x e^y$ 
\end{enumerate}

% beg ... i





\item Let $a$ be a non-zero constant.  Suppose the solution $y(x)$ to the Bernoulli IVP $\displaystyle y^{\,\prime}=ay+\frac{4xe^{3ax}}{y^2}$, $y(0)=1$ satisfies $y(1)=8$.  Find $a$ to the nearest hundredth. 

    %   y^{\,\prime}-ay=4xe^{3ax}y^{-2} so y_1=e^{ax}.  So y=ue^{ax} and hence u'e^{ax}=4xe^{3ax} (ue^{ax})^{-2} =4xe^{3ax} u^{-2} e^{-2ax} so u^2 u'=4x.

    %  So u^3 /3 = 2x^2+C so u =(6x^2+C)^{1/3} and y=e^{ax} (6x^2+C)^{1/3}.  Since y(0)=1, C=1.  So y=(6x^2+1)^{1/3} e^{ax}

    %  Check: y^{\,\prime}=4x(6x^2+1)^{-2/3} e^{ax}+a(6x^2+1)^{1/3} e^{ax}=4xe^{3ax}y^{-2} +ay checks out

    %  8=(7)^{1/3} e^{a} so ln (8/ 7^(1/3))\approx 1.43


\newpage

\item Determine which of the following equations are ``nonlinear homogeneous."  Select all that apply.

\begin{enumerate}[label=(\alph*)]
\item $\displaystyle y\,'=1+\frac{y}{x}$
\item $\displaystyle y\,'=1+\frac{x}{x+y}$ % ***
\item $\displaystyle y\,'=\frac{2x}{x^2+3y}$ 
\item $\displaystyle y\,'=\frac{2x^2}{x^2+y^2}$ % ***
\item $\displaystyle y\,'=e^{x/(x+y)}$  % ****
\item $\displaystyle y\,'=x+\frac{x}{y}+\frac{x^2}{y}$
\end{enumerate}

% (b), (d), (e) with options thru (f)


\item Given the nonlinear homogeneous differential equation $\displaystyle y^{\,\prime}=\frac{xy}{x^2-y^2}$ determine the transformed separable equation (in standard form) that is the result of applying the substitution $u=\frac{y}{x}$.

\begin{enumerate}[label=(\alph*)]
\item $\displaystyle \frac{1-u^2}{u^3} u'=\frac{1}{x}$ % ***
\item $\displaystyle \frac{1-u^3}{u^2} u'=\frac{1}{x}$
\item $\displaystyle \frac{1}{1-u^2} u'=-\frac{1}{x}$
\item $\displaystyle \frac{u^3}{1-u^2} u'=\frac{1}{x}$
\item $\displaystyle \frac{u^2}{1-u^3} u'=\frac{1}{x}$
\item None of these
\end{enumerate}


% (a) ... (f)


% u' x + u=\frac{u}{1-u^2} so u' x=-u+\frac{u}{1-u^2}=\frac{u^3}{1-u^2}

% \frac{1-u^2}{u^3} u'=\frac{1}{x}

\item  Suppose $y(x)$ satisfies the nonlinear homogeneous IVP $\displaystyle y^{\,\prime}=\frac{x^2+y^2}{xy}$, $y(1)=1$.  Find $y(2026)$.  Round to the nearest whole number.

% u' x+u=u+1/u so u'x=1/u so u u'=1/x so 0.5u^2 =ln(x)+C so u^2=ln(x^2)+C.  So (y/x)^2=ln(x^2)+C.  C=1.  So y=\sqrt{x^2 ln(x^2)+x^2}.  y(2026)=2026*Sqrt(2*ln(2026)+1)\approx 8161

\newpage

\item Suppose $y(x)$ satisfies $y\,'=-\sqrt{x+y}$, $y(0)=0$.  Which of the following is an implicit solution for $y(x)$?

% u=x+y.  y^{\,\prime}=u'-1=-sqrt(u) ... so u'=1-sqrt(u)

% 1/(1-sqrt(u)) u'=1 ---->

% Integral of 1/(1-sqrt(u)) du:  Let v=1-sqrt(u). Then (v-1)^2=u, 2(v-1)dv=du.  So Integral of 2(v-1)/v dv=\int 2-2/v dv=2-2ln(v)+C

% So y=2-2/(1-sqrt(x+y))-x

\begin{enumerate}[label=(\alph*)]
\item $\sqrt{x+y}+\ln(1-\sqrt{x+y}\,)=-\dfrac{x}{2}$  % ***
\item $\sqrt{x+y}-\ln(1-\sqrt{x+y}\,)=-\dfrac{x}{2}$
\item $\sqrt{x+y}+\ln(1-\sqrt{x+y}\,)=-x$
\item $\sqrt{x+y}-\ln(1-\sqrt{x+y}\,)=-x$
\item $y - \ln|x+y+1|=-\dfrac{x}{2}$
\item $y+\ln|x+y+1|=-\dfrac{x}{2}$
\end{enumerate}

% (a) with options thru (f)


\item Suppose $y(x)$ satisfies $y^{\,\prime}=(9x-y+1)^{2}$, $y(0)=0$.  Find $y(1)$.  Round to the nearest hundredth.

% u=9x-y+1. u'=9-y^{\,\prime} so y^{\,\prime}=9-u'=u^2 implies u'=9-u^2 and 1/[(3-u)(3+u)] u'=1.

% Then  [ 1/(3-u) + 1/(3+u) ] u'=6 implies ln[(3+u)/(3-u)]=6x+C implies (3+u)/(3-u)=Ce^(6x).  -1 + 6/(3-u)=Ce^(4x)

% So 6/(3-u)=1+Ce^(6x)

% So 3-u=6/(1+Ce^(6x))

% So 3-(9x-y+1)=6/(1+Ce^(6x))

% So y=9x-2 + 6/(1+Ce^(6x))

% C=2.  y=9x-2 +6/(1+2e^(6x)).  Then y(1)=9(1)-2 +6/(1+2e^(6))\approx 7.01



% Check: y^{\,\prime}=9 - 72e^(6x)/(1+2e^(6x))^2=9[(1+2e^(6x))^2 -8e^(6x)]/ (1+2e^(6x))^2=9(1-2e^(6x))^2/ (1+2e^(6x))^2.

% (9x-y+1)^{2}=[3-6/(1+2e^(6x))]^2=9[1-2/(1+2e^(6x))]^2=9[((2e^(6x))-1)/(1+2e^(6x))]^2 checks out.

\item Solve the IVP $xy^{\,\prime}=y+\sqrt{x^2+y^2}$, $y(1)=0$.  Determine $x$ such that $y(x)=2$.  Round to the nearest hundredth.

% 2.24

Hint: you may find $\displaystyle \int \sec(\theta)\,\mathrm{d}\theta=\ln{|\sec (\theta)+\tan (\theta)|}+C$ useful.

% y' = y/x + sqrt(1+(y/x)^2)

% y=ux --> y'=u'x+u=u+sqrt(1+u^2)

% u'x = sqrt(1+u^2)

% 1/sqrt(1+u^2) u' = 1/x

% u=tan(t).  u'=sec^2(t).  sqrt(1+u^2)=sec(t)

% sec(t) dt = 1/x dx

% ln|sec(t)+tan(t)| = ln|x|+C 

% sqrt(1+u^2)+u = Cx 

% sqrt(1+(y/x))^2)+y/x = Cx 

% 1 = C 

% sqrt(x^2+y^2)+y = x^2 

% sqrt(x^2+4)+2 = x^2 

% sqrt(x^2+4) = x^2 -2

% x^2+4 = x^4 -4x^2+4

% 0 = x^4 -5x^2 ... x=sqrt(5)\approx 2.24


\end{enumerate}

\end{document}
